\documentclass[11pt]{article}
\usepackage{amsmath,amssymb,hyperref}
\title{Finite Joint-Family Certificates with Exact Axis Cancellation}
\author{Research checkpoint, branch ar/027ac8030b43}
\date{November 24, 2026}
\begin{document}
\maketitle

\paragraph{Status.}
These are finite computations checked by exact rational arithmetic.
They are not Lean declarations, a universal decoder theorem, or an
asymptotic upper bound. Optimizer failures are not impossibility results.
The original head model and worst-case complexity are unchanged.

\section{The tested interface}

For $x\in\{0,1\}^k$ and a one-hot column $j\in[s]$, the score is
\[
F_j(x)=\tau_j+\sum_{h=1}^{H}\left(
\frac{A_h(x)+a_{hj}}{b_{hj}+d_h\cdot x}
-\frac{A_h(x)}{B+d_h\cdot x}\right),\qquad
A_h(x)=\sum_{i=1}^k\alpha_{hi}x_i.
\tag{1}
\]
Every $d_{hi}>0$, and $1\leq b_{hj}\leq1+\delta$. The shared directions
$d_h$ and row numerators $A_h$ may depend jointly on the entire family of
column labels. They are not chosen separately for each column.

To embed the pairs on the full Boolean cube, choose a common intercept
$E>\max(B,1+\delta)$. Their denominators at arbitrary column bits $w$ are
\[
E+d_h\cdot x-\sum_j(E-b_{hj})w_j,
\qquad
E+d_h\cdot x-\sum_j(E-B)w_j.
\tag{2}
\]
The first numerator is $A_h(x)+\sum_j a_{hj}w_j$, and the second is
$-A_h(x)$. At $w=0$ their denominators agree and their numerators cancel
exactly. Complementing all row bits converts the row slopes in (2) to
strictly negative slopes, matching the reverse compiler orientation.
All column slopes are already strictly negative. Full-cube positivity
follows from the exact inequalities
\[
E-\sum_j(E-b_{hj})>0,\qquad E-s(E-B)>0.
\tag{3}
\]
These conditions are checked separately for every certificate, not inferred
from positivity only at the tested one-hot inputs. The column biases
$\tau_j$ require the separate admissible bias realization used by the
transfer. Thus the advertised local budget is $2H+1$ heads, not $H$.

\section{Results and reproduction}

The committed file \texttt{LargeDomainJointFamilyCertificates.json}
contains 20 certified families with
\[
k=5,\quad s=5,\quad H=6,\quad
\delta=\frac1{20},\quad B=\frac{11}{10},\quad E=\frac{23}{20}.
\]
The first target includes coordinate and full-parity columns. The other
19 targets have independently sampled signs at every row and column.
Each family specifies all rational parameters in (1), all 160 target
signs, and its exact positive minimum margin. All 20 families pass.
This gives 3200 strict signs in total, checked without a tolerance.
The independently checked full-cube minimum denominator lower bound
in (3) is positive for both members of every pair.

Run the verifier from the repository root:
\begin{verbatim}
python3 artifacts/head-compression/LargeDomainJointFamilyVerify.py \
  artifacts/head-compression/LargeDomainJointFamilyCertificates.json
\end{verbatim}
The verifier uses only Python integers and \texttt{fractions.Fraction};
it does not import the search, a numerical linear algebra library, or an
optimization package. It reconstructs every score, checks its sign,
compares the exact minimum margin with the stored value, and checks (3).
The general algebraic full-cube lift is described above; the global
cover and preservation of previously assigned signs remain separate
proof obligations handled by the transfer theorem.

The generating script is \texttt{LargeDomainJointFamilySearch.py}. It
uses variable projection for the linear numerator coefficients and
optimizes the positive directions and compact offsets. Positive target
magnitudes are also optimized, so the search asks for membership in a
sign cone rather than exact fitting of the values $\pm1$. A numerical
linear program proposes final coefficients. Decimal rationalization is
then checked with exact arithmetic. The shared machine command is
prefixed by \texttt{heavy}.

The committed certificates were generated with the equivalent of:
\begin{verbatim}
heavy python3 artifacts/head-compression/LargeDomainJointFamilySearch.py \
  --dimension 5 --heads 6 --columns 5 --targets 20 --restarts 20 \
  --iterations 500 --affine-row --paired --sign-cone \
  --baseline 1.1 --delta .05 --seed 241124 --output result.json
\end{verbatim}
The explicit baseline option was added immediately after this run; its
original value was $1+2\delta=1.1$. This provenance is recorded in the
certificate file. Subsequent searches can use baseline one, the form of
the general paired compact-family transfer.

\section{Limits of the evidence}

These certificates show that sharing directions and row coefficients
across several unrelated columns is feasible at this small size.
They do not show that every five-column family has such a representation.
Nor do they show $H(k)=o(2^k/\log k)$, which would be needed to improve the
inherited asymptotic scale through a full-row-cube transfer.

An earlier unpaired $k=4,H=3$ experiment failed on a column of full
four-bit parity. That failure has an independent mathematical explanation:
a sum of three affine ratios with positive denominators, plus a bias,
has after clearing denominators an input polynomial of degree at most
three; strict four-bit parity has threshold degree four. This observation
does not apply to the paired six-head family in (1), and is not used as
an interpretation of failures for the larger paired experiments.
\end{document}
